\documentclass{article}

\usepackage[T1]{fontenc}
\usepackage{iclr2027_conference,times}
%%%%% NEW MATH DEFINITIONS %%%%%

\usepackage{amsmath,amsfonts,bm}

% Mark sections of captions for referring to divisions of figures
\newcommand{\figleft}{{\em (Left)}}
\newcommand{\figcenter}{{\em (Center)}}
\newcommand{\figright}{{\em (Right)}}
\newcommand{\figtop}{{\em (Top)}}
\newcommand{\figbottom}{{\em (Bottom)}}
\newcommand{\captiona}{{\em (a)}}
\newcommand{\captionb}{{\em (b)}}
\newcommand{\captionc}{{\em (c)}}
\newcommand{\captiond}{{\em (d)}}

% Highlight a newly defined term
\newcommand{\newterm}[1]{{\bf #1}}


% Figure reference, lower-case.
\def\figref#1{figure~\ref{#1}}
% Figure reference, capital. For start of sentence
\def\Figref#1{Figure~\ref{#1}}
\def\twofigref#1#2{figures \ref{#1} and \ref{#2}}
\def\quadfigref#1#2#3#4{figures \ref{#1}, \ref{#2}, \ref{#3} and \ref{#4}}
% Section reference, lower-case.
\def\secref#1{section~\ref{#1}}
% Section reference, capital.
\def\Secref#1{Section~\ref{#1}}
% Reference to two sections.
\def\twosecrefs#1#2{sections \ref{#1} and \ref{#2}}
% Reference to three sections.
\def\secrefs#1#2#3{sections \ref{#1}, \ref{#2} and \ref{#3}}
% Reference to an equation, lower-case.
\def\eqref#1{equation~\ref{#1}}
% Reference to an equation, upper case
\def\Eqref#1{Equation~\ref{#1}}
% A raw reference to an equation---avoid using if possible
\def\plaineqref#1{\ref{#1}}
% Reference to a chapter, lower-case.
\def\chapref#1{chapter~\ref{#1}}
% Reference to an equation, upper case.
\def\Chapref#1{Chapter~\ref{#1}}
% Reference to a range of chapters
\def\rangechapref#1#2{chapters\ref{#1}--\ref{#2}}
% Reference to an algorithm, lower-case.
\def\algref#1{algorithm~\ref{#1}}
% Reference to an algorithm, upper case.
\def\Algref#1{Algorithm~\ref{#1}}
\def\twoalgref#1#2{algorithms \ref{#1} and \ref{#2}}
\def\Twoalgref#1#2{Algorithms \ref{#1} and \ref{#2}}
% Reference to a part, lower case
\def\partref#1{part~\ref{#1}}
% Reference to a part, upper case
\def\Partref#1{Part~\ref{#1}}
\def\twopartref#1#2{parts \ref{#1} and \ref{#2}}

\def\ceil#1{\lceil #1 \rceil}
\def\floor#1{\lfloor #1 \rfloor}
\def\1{\bm{1}}
\newcommand{\train}{\mathcal{D}}
\newcommand{\valid}{\mathcal{D_{\mathrm{valid}}}}
\newcommand{\test}{\mathcal{D_{\mathrm{test}}}}

\def\eps{{\epsilon}}


% Random variables
\def\reta{{\textnormal{$\eta$}}}
\def\ra{{\textnormal{a}}}
\def\rb{{\textnormal{b}}}
\def\rc{{\textnormal{c}}}
\def\rd{{\textnormal{d}}}
\def\re{{\textnormal{e}}}
\def\rf{{\textnormal{f}}}
\def\rg{{\textnormal{g}}}
\def\rh{{\textnormal{h}}}
\def\ri{{\textnormal{i}}}
\def\rj{{\textnormal{j}}}
\def\rk{{\textnormal{k}}}
\def\rl{{\textnormal{l}}}
% rm is already a command, just don't name any random variables m
\def\rn{{\textnormal{n}}}
\def\ro{{\textnormal{o}}}
\def\rp{{\textnormal{p}}}
\def\rq{{\textnormal{q}}}
\def\rr{{\textnormal{r}}}
\def\rs{{\textnormal{s}}}
\def\rt{{\textnormal{t}}}
\def\ru{{\textnormal{u}}}
\def\rv{{\textnormal{v}}}
\def\rw{{\textnormal{w}}}
\def\rx{{\textnormal{x}}}
\def\ry{{\textnormal{y}}}
\def\rz{{\textnormal{z}}}

% Random vectors
\def\rvepsilon{{\mathbf{\epsilon}}}
\def\rvtheta{{\mathbf{\theta}}}
\def\rva{{\mathbf{a}}}
\def\rvb{{\mathbf{b}}}
\def\rvc{{\mathbf{c}}}
\def\rvd{{\mathbf{d}}}
\def\rve{{\mathbf{e}}}
\def\rvf{{\mathbf{f}}}
\def\rvg{{\mathbf{g}}}
\def\rvh{{\mathbf{h}}}
\def\rvu{{\mathbf{i}}}
\def\rvj{{\mathbf{j}}}
\def\rvk{{\mathbf{k}}}
\def\rvl{{\mathbf{l}}}
\def\rvm{{\mathbf{m}}}
\def\rvn{{\mathbf{n}}}
\def\rvo{{\mathbf{o}}}
\def\rvp{{\mathbf{p}}}
\def\rvq{{\mathbf{q}}}
\def\rvr{{\mathbf{r}}}
\def\rvs{{\mathbf{s}}}
\def\rvt{{\mathbf{t}}}
\def\rvu{{\mathbf{u}}}
\def\rvv{{\mathbf{v}}}
\def\rvw{{\mathbf{w}}}
\def\rvx{{\mathbf{x}}}
\def\rvy{{\mathbf{y}}}
\def\rvz{{\mathbf{z}}}

% Elements of random vectors
\def\erva{{\textnormal{a}}}
\def\ervb{{\textnormal{b}}}
\def\ervc{{\textnormal{c}}}
\def\ervd{{\textnormal{d}}}
\def\erve{{\textnormal{e}}}
\def\ervf{{\textnormal{f}}}
\def\ervg{{\textnormal{g}}}
\def\ervh{{\textnormal{h}}}
\def\ervi{{\textnormal{i}}}
\def\ervj{{\textnormal{j}}}
\def\ervk{{\textnormal{k}}}
\def\ervl{{\textnormal{l}}}
\def\ervm{{\textnormal{m}}}
\def\ervn{{\textnormal{n}}}
\def\ervo{{\textnormal{o}}}
\def\ervp{{\textnormal{p}}}
\def\ervq{{\textnormal{q}}}
\def\ervr{{\textnormal{r}}}
\def\ervs{{\textnormal{s}}}
\def\ervt{{\textnormal{t}}}
\def\ervu{{\textnormal{u}}}
\def\ervv{{\textnormal{v}}}
\def\ervw{{\textnormal{w}}}
\def\ervx{{\textnormal{x}}}
\def\ervy{{\textnormal{y}}}
\def\ervz{{\textnormal{z}}}

% Random matrices
\def\rmA{{\mathbf{A}}}
\def\rmB{{\mathbf{B}}}
\def\rmC{{\mathbf{C}}}
\def\rmD{{\mathbf{D}}}
\def\rmE{{\mathbf{E}}}
\def\rmF{{\mathbf{F}}}
\def\rmG{{\mathbf{G}}}
\def\rmH{{\mathbf{H}}}
\def\rmI{{\mathbf{I}}}
\def\rmJ{{\mathbf{J}}}
\def\rmK{{\mathbf{K}}}
\def\rmL{{\mathbf{L}}}
\def\rmM{{\mathbf{M}}}
\def\rmN{{\mathbf{N}}}
\def\rmO{{\mathbf{O}}}
\def\rmP{{\mathbf{P}}}
\def\rmQ{{\mathbf{Q}}}
\def\rmR{{\mathbf{R}}}
\def\rmS{{\mathbf{S}}}
\def\rmT{{\mathbf{T}}}
\def\rmU{{\mathbf{U}}}
\def\rmV{{\mathbf{V}}}
\def\rmW{{\mathbf{W}}}
\def\rmX{{\mathbf{X}}}
\def\rmY{{\mathbf{Y}}}
\def\rmZ{{\mathbf{Z}}}

% Elements of random matrices
\def\ermA{{\textnormal{A}}}
\def\ermB{{\textnormal{B}}}
\def\ermC{{\textnormal{C}}}
\def\ermD{{\textnormal{D}}}
\def\ermE{{\textnormal{E}}}
\def\ermF{{\textnormal{F}}}
\def\ermG{{\textnormal{G}}}
\def\ermH{{\textnormal{H}}}
\def\ermI{{\textnormal{I}}}
\def\ermJ{{\textnormal{J}}}
\def\ermK{{\textnormal{K}}}
\def\ermL{{\textnormal{L}}}
\def\ermM{{\textnormal{M}}}
\def\ermN{{\textnormal{N}}}
\def\ermO{{\textnormal{O}}}
\def\ermP{{\textnormal{P}}}
\def\ermQ{{\textnormal{Q}}}
\def\ermR{{\textnormal{R}}}
\def\ermS{{\textnormal{S}}}
\def\ermT{{\textnormal{T}}}
\def\ermU{{\textnormal{U}}}
\def\ermV{{\textnormal{V}}}
\def\ermW{{\textnormal{W}}}
\def\ermX{{\textnormal{X}}}
\def\ermY{{\textnormal{Y}}}
\def\ermZ{{\textnormal{Z}}}

% Vectors
\def\vzero{{\bm{0}}}
\def\vone{{\bm{1}}}
\def\vmu{{\bm{\mu}}}
\def\vtheta{{\bm{\theta}}}
\def\va{{\bm{a}}}
\def\vb{{\bm{b}}}
\def\vc{{\bm{c}}}
\def\vd{{\bm{d}}}
\def\ve{{\bm{e}}}
\def\vf{{\bm{f}}}
\def\vg{{\bm{g}}}
\def\vh{{\bm{h}}}
\def\vi{{\bm{i}}}
\def\vj{{\bm{j}}}
\def\vk{{\bm{k}}}
\def\vl{{\bm{l}}}
\def\vm{{\bm{m}}}
\def\vn{{\bm{n}}}
\def\vo{{\bm{o}}}
\def\vp{{\bm{p}}}
\def\vq{{\bm{q}}}
\def\vr{{\bm{r}}}
\def\vs{{\bm{s}}}
\def\vt{{\bm{t}}}
\def\vu{{\bm{u}}}
\def\vv{{\bm{v}}}
\def\vw{{\bm{w}}}
\def\vx{{\bm{x}}}
\def\vy{{\bm{y}}}
\def\vz{{\bm{z}}}

% Elements of vectors
\def\evalpha{{\alpha}}
\def\evbeta{{\beta}}
\def\evepsilon{{\epsilon}}
\def\evlambda{{\lambda}}
\def\evomega{{\omega}}
\def\evmu{{\mu}}
\def\evpsi{{\psi}}
\def\evsigma{{\sigma}}
\def\evtheta{{\theta}}
\def\eva{{a}}
\def\evb{{b}}
\def\evc{{c}}
\def\evd{{d}}
\def\eve{{e}}
\def\evf{{f}}
\def\evg{{g}}
\def\evh{{h}}
\def\evi{{i}}
\def\evj{{j}}
\def\evk{{k}}
\def\evl{{l}}
\def\evm{{m}}
\def\evn{{n}}
\def\evo{{o}}
\def\evp{{p}}
\def\evq{{q}}
\def\evr{{r}}
\def\evs{{s}}
\def\evt{{t}}
\def\evu{{u}}
\def\evv{{v}}
\def\evw{{w}}
\def\evx{{x}}
\def\evy{{y}}
\def\evz{{z}}

% Matrix
\def\mA{{\bm{A}}}
\def\mB{{\bm{B}}}
\def\mC{{\bm{C}}}
\def\mD{{\bm{D}}}
\def\mE{{\bm{E}}}
\def\mF{{\bm{F}}}
\def\mG{{\bm{G}}}
\def\mH{{\bm{H}}}
\def\mI{{\bm{I}}}
\def\mJ{{\bm{J}}}
\def\mK{{\bm{K}}}
\def\mL{{\bm{L}}}
\def\mM{{\bm{M}}}
\def\mN{{\bm{N}}}
\def\mO{{\bm{O}}}
\def\mP{{\bm{P}}}
\def\mQ{{\bm{Q}}}
\def\mR{{\bm{R}}}
\def\mS{{\bm{S}}}
\def\mT{{\bm{T}}}
\def\mU{{\bm{U}}}
\def\mV{{\bm{V}}}
\def\mW{{\bm{W}}}
\def\mX{{\bm{X}}}
\def\mY{{\bm{Y}}}
\def\mZ{{\bm{Z}}}
\def\mBeta{{\bm{\beta}}}
\def\mPhi{{\bm{\Phi}}}
\def\mLambda{{\bm{\Lambda}}}
\def\mSigma{{\bm{\Sigma}}}

% Tensor
\DeclareMathAlphabet{\mathsfit}{\encodingdefault}{\sfdefault}{m}{sl}
\SetMathAlphabet{\mathsfit}{bold}{\encodingdefault}{\sfdefault}{bx}{n}
\newcommand{\tens}[1]{\bm{\mathsfit{#1}}}
\def\tA{{\tens{A}}}
\def\tB{{\tens{B}}}
\def\tC{{\tens{C}}}
\def\tD{{\tens{D}}}
\def\tE{{\tens{E}}}
\def\tF{{\tens{F}}}
\def\tG{{\tens{G}}}
\def\tH{{\tens{H}}}
\def\tI{{\tens{I}}}
\def\tJ{{\tens{J}}}
\def\tK{{\tens{K}}}
\def\tL{{\tens{L}}}
\def\tM{{\tens{M}}}
\def\tN{{\tens{N}}}
\def\tO{{\tens{O}}}
\def\tP{{\tens{P}}}
\def\tQ{{\tens{Q}}}
\def\tR{{\tens{R}}}
\def\tS{{\tens{S}}}
\def\tT{{\tens{T}}}
\def\tU{{\tens{U}}}
\def\tV{{\tens{V}}}
\def\tW{{\tens{W}}}
\def\tX{{\tens{X}}}
\def\tY{{\tens{Y}}}
\def\tZ{{\tens{Z}}}


% Graph
\def\gA{{\mathcal{A}}}
\def\gB{{\mathcal{B}}}
\def\gC{{\mathcal{C}}}
\def\gD{{\mathcal{D}}}
\def\gE{{\mathcal{E}}}
\def\gF{{\mathcal{F}}}
\def\gG{{\mathcal{G}}}
\def\gH{{\mathcal{H}}}
\def\gI{{\mathcal{I}}}
\def\gJ{{\mathcal{J}}}
\def\gK{{\mathcal{K}}}
\def\gL{{\mathcal{L}}}
\def\gM{{\mathcal{M}}}
\def\gN{{\mathcal{N}}}
\def\gO{{\mathcal{O}}}
\def\gP{{\mathcal{P}}}
\def\gQ{{\mathcal{Q}}}
\def\gR{{\mathcal{R}}}
\def\gS{{\mathcal{S}}}
\def\gT{{\mathcal{T}}}
\def\gU{{\mathcal{U}}}
\def\gV{{\mathcal{V}}}
\def\gW{{\mathcal{W}}}
\def\gX{{\mathcal{X}}}
\def\gY{{\mathcal{Y}}}
\def\gZ{{\mathcal{Z}}}

% Sets
\def\sA{{\mathbb{A}}}
\def\sB{{\mathbb{B}}}
\def\sC{{\mathbb{C}}}
\def\sD{{\mathbb{D}}}
% Don't use a set called E, because this would be the same as our symbol
% for expectation.
\def\sF{{\mathbb{F}}}
\def\sG{{\mathbb{G}}}
\def\sH{{\mathbb{H}}}
\def\sI{{\mathbb{I}}}
\def\sJ{{\mathbb{J}}}
\def\sK{{\mathbb{K}}}
\def\sL{{\mathbb{L}}}
\def\sM{{\mathbb{M}}}
\def\sN{{\mathbb{N}}}
\def\sO{{\mathbb{O}}}
\def\sP{{\mathbb{P}}}
\def\sQ{{\mathbb{Q}}}
\def\sR{{\mathbb{R}}}
\def\sS{{\mathbb{S}}}
\def\sT{{\mathbb{T}}}
\def\sU{{\mathbb{U}}}
\def\sV{{\mathbb{V}}}
\def\sW{{\mathbb{W}}}
\def\sX{{\mathbb{X}}}
\def\sY{{\mathbb{Y}}}
\def\sZ{{\mathbb{Z}}}

% Entries of a matrix
\def\emLambda{{\Lambda}}
\def\emA{{A}}
\def\emB{{B}}
\def\emC{{C}}
\def\emD{{D}}
\def\emE{{E}}
\def\emF{{F}}
\def\emG{{G}}
\def\emH{{H}}
\def\emI{{I}}
\def\emJ{{J}}
\def\emK{{K}}
\def\emL{{L}}
\def\emM{{M}}
\def\emN{{N}}
\def\emO{{O}}
\def\emP{{P}}
\def\emQ{{Q}}
\def\emR{{R}}
\def\emS{{S}}
\def\emT{{T}}
\def\emU{{U}}
\def\emV{{V}}
\def\emW{{W}}
\def\emX{{X}}
\def\emY{{Y}}
\def\emZ{{Z}}
\def\emSigma{{\Sigma}}

% entries of a tensor
% Same font as tensor, without \bm wrapper
\newcommand{\etens}[1]{\mathsfit{#1}}
\def\etLambda{{\etens{\Lambda}}}
\def\etA{{\etens{A}}}
\def\etB{{\etens{B}}}
\def\etC{{\etens{C}}}
\def\etD{{\etens{D}}}
\def\etE{{\etens{E}}}
\def\etF{{\etens{F}}}
\def\etG{{\etens{G}}}
\def\etH{{\etens{H}}}
\def\etI{{\etens{I}}}
\def\etJ{{\etens{J}}}
\def\etK{{\etens{K}}}
\def\etL{{\etens{L}}}
\def\etM{{\etens{M}}}
\def\etN{{\etens{N}}}
\def\etO{{\etens{O}}}
\def\etP{{\etens{P}}}
\def\etQ{{\etens{Q}}}
\def\etR{{\etens{R}}}
\def\etS{{\etens{S}}}
\def\etT{{\etens{T}}}
\def\etU{{\etens{U}}}
\def\etV{{\etens{V}}}
\def\etW{{\etens{W}}}
\def\etX{{\etens{X}}}
\def\etY{{\etens{Y}}}
\def\etZ{{\etens{Z}}}

% The true underlying data generating distribution
\newcommand{\pdata}{p_{\rm{data}}}
% The empirical distribution defined by the training set
\newcommand{\ptrain}{\hat{p}_{\rm{data}}}
\newcommand{\Ptrain}{\hat{P}_{\rm{data}}}
% The model distribution
\newcommand{\pmodel}{p_{\rm{model}}}
\newcommand{\Pmodel}{P_{\rm{model}}}
\newcommand{\ptildemodel}{\tilde{p}_{\rm{model}}}
% Stochastic autoencoder distributions
\newcommand{\pencode}{p_{\rm{encoder}}}
\newcommand{\pdecode}{p_{\rm{decoder}}}
\newcommand{\precons}{p_{\rm{reconstruct}}}

\newcommand{\laplace}{\mathrm{Laplace}} % Laplace distribution

\newcommand{\E}{\mathbb{E}}
\newcommand{\Ls}{\mathcal{L}}
\newcommand{\R}{\mathbb{R}}
\newcommand{\emp}{\tilde{p}}
\newcommand{\lr}{\alpha}
\newcommand{\reg}{\lambda}
\newcommand{\rect}{\mathrm{rectifier}}
\newcommand{\softmax}{\mathrm{softmax}}
\newcommand{\sigmoid}{\sigma}
\newcommand{\softplus}{\zeta}
\newcommand{\KL}{D_{\mathrm{KL}}}
\newcommand{\Var}{\mathrm{Var}}
\newcommand{\standarderror}{\mathrm{SE}}
\newcommand{\Cov}{\mathrm{Cov}}
% Wolfram Mathworld says $L^2$ is for function spaces and $\ell^2$ is for vectors
% But then they seem to use $L^2$ for vectors throughout the site, and so does
% wikipedia.
\newcommand{\normlzero}{L^0}
\newcommand{\normlone}{L^1}
\newcommand{\normltwo}{L^2}
\newcommand{\normlp}{L^p}
\newcommand{\normmax}{L^\infty}

\newcommand{\parents}{Pa} % See usage in notation.tex. Chosen to match Daphne's book.

\DeclareMathOperator*{\argmax}{arg\,max}
\DeclareMathOperator*{\argmin}{arg\,min}

\DeclareMathOperator{\sign}{sign}
\DeclareMathOperator{\Tr}{Tr}
\let\ab\allowbreak


\usepackage{float}
\usepackage{hyperref}
\usepackage{url}
\usepackage{graphicx}
\usepackage{amsmath,amssymb,bm}
\usepackage{booktabs,array,multirow,threeparttable,adjustbox}
\usepackage{placeins}
\usepackage{tikz}
\usetikzlibrary{arrows.meta,positioning,fit,calc}
\hypersetup{hidelinks}

\title{RAL-MoE-ROM: Hierarchical Mixture-of-Experts for Regime-Adaptive Reduced-Order Model }

% Anonymous submission. Restore the author block only for the camera-ready version.
\author{Anonymous Authors}

% Keep disabled for the anonymous submission.
% \iclrfinalcopy


\begin{document}
\raggedbottom

\maketitle

\begin{abstract}
Reduced-order models (ROMs) approximate high-dimensional dynamical systems using a small number of state variables. However, their predictive accuracy can deteriorate when parameter changes lead to different dynamical regimes. Different regimes may require different reduced representations and dynamics, making it difficult for a single model to remain accurate across the full parameter range. We introduce RAL-MoE-ROM, a hierarchical mixture-of-experts framework with two levels: regime-local specialists and cross-regime routing with fusion.. At the first level, each local model, called a specialist, uses a structured sparse mixture of experts to improve its reduced dynamics and is trained through multi-step autonomous rollouts. At the second level, an outer router selects a specialist or a pair of neighboring specialists. When two specialists are selected, a learned gate uses the initial physical history to combine their reconstructed fields. Across three flow benchmarks spanning steady, Hopf-onset, and periodic regimes, the local specialists provide accurate autonomous predictions, while cross-regime routing further improve prediction in the evaluated regime overlaps.
\end{abstract}

\section{Introduction}

High-fidelity simulations of high-dimensional dynamical systems
can be expensive when repeated predictions are required across many parameter settings. Reduced-order models (ROMs) provide a cost-effective surrogate by representing the physical state with a small number of state variables. The problem becomes more challenging when changing a parameter also changes the type of regimes. In incompressible flows, for example, varying the Reynolds number can lead from relaxation toward a steady state to oscillatory growth and sustained periodic motion. A useful predictive model must capture these regimes while evolving autonomously.



Such regime variation introduces two closely related challenges for ROMs. The first concerns representation: spatial structures that are important in one regime may differ from those required in another, making a single representation less effective. The second is dynamical: the reduced evolution law must accommodate qualitatively different behaviors, including stable relaxation, instability growth, and nonlinear periodic oscillations. Increasing model capacity alone cannot compensate for an inadequate reduced representation, while improving the representation does not by itself guarantee accurate autonomous rollout. These observations motivate adapting the local representation and the local dynamics separately, while providing a consistent way to assemble their predictions.




Projection-based ROMs provide an explicit reduced representation and reduced dynamics, commonly through proper orthogonal decomposition (POD) and Galerkin projection \citep{sirovich1987,berkooz1993,benner2015}. However,projection-based ROMs are inherently affected by truncation, which neglects the influence of discarded modes on the retained coordinates, leading to closure errors that degrade long-horizon prediction.\citep{amsallem2012}. Data-driven ROMs offer greater flexibility in approximating nonlinear dynamics, but low one-step error does not guarantee accurate autonomous rollouts, and they often lack physical interpretability\citep{duraisamy2019,brunton2020}. Hybrid ROMs combine projection-based ROMs with learned corrections or learned reduced representations, retaining physical structure while increasing approximation flexibility
\citep{dar2023,lee2020,fresca2022,geelen2023}. These developments improve the modeling of reduced dynamics, but they do not by themselves resolve the representational difficulty that arises when qualitatively different regimes must be described over a broad parameter range.





A complementary strategy localizes the reduced representation, using multiple subspaces, clusters, or manifolds for different portions of the solution set \citep{kaiser2014,li2021cluster,diaz2024}. These directions address complementary aspects of the broad-parameter problem: hybrid modeling improves the fidelity of the reduced dynamics, whereas localization adapts the reduced representation to regime-dependent structures. Localization, however, creates a new assembly problem.  Independently constructed local ROMs generally have different reduced state variables and  operators, so their reduced coordinates are chart-specific. Direct interpolation or averaging in reduced-coordinate space therefore has no natural cross-chart interpretation. Mixture-of-experts (MoE) architectures provide a natural framework for conditional specialization \citep{jacobs1991,jordan1994,shazeer2017,fedus2022}, but standard MoE formulations typically route specialized predictors within a shared representation or common feature space. Here, in contrast, the experts are complete local ROMs whose coordinates are not shared. The problem is therefore not only routing among local ROMs, but also assembling or fusing their outputs across representations. This makes the problem hierarchical: one must adapt the reduced dynamics within each local representation and, separately, route to and fuse complete local ROMs across representations. This raises a distinct question that standard within-representation MoE routing does not address: how can independently constructed local dynamical models be routed and assembled without forcing them into a common latent coordinate system?


To address this gap, we introduce the Regime-Adaptive Localized Mixture-of-Experts Reduced-Order Model (RAL-MoE-ROM), a two-level hierarchical framework that separates regime-local reduced dynamics from cross-regime routing and fusion. The parameter domain is covered by prescribed, overlapping regime-local regions, each equipped with an independent  autonomous specialist. At the first level, each specialist, use an explicit projected-based ROM  and augments it with sparsely  structured MOEs which trained trained through  multi-step autonomous rollouts. Rather than using conventional MLP-only experts, each activated expert decomposes its output into three additive components: a nonlinear neural branch, together with explicit linear and low-rank quadratic branches derived from the projection-based. At the second level, a parameter-conditioned outer router produces window-level routing weights over the local specialists, while a fixed applicability mask determines whether inference uses one specialist or a pair of adjacent specialists. When two specialists are selected, a learned fusion gate combines their reconstructed physical fields conditioned on the initial physical history. The overall architecture is illustrated in Figure~\ref{fig:ral_overview}, highlighting the regime-local specialists and the cross-regime routing and fusion mechanism.




Our main contributions are as follows:
\begin{itemize}

\item \textbf{Structured regime-local specialists for autonomous reduced dynamics.} Each specialist combines a physics-based ROM with a structured sparse MoE closure and is trained under closed-loop multi-step rollouts.

\item \textbf{Cross-regime routing and fusion.}
An outer router selects a specialist or a pair of neighboring specialists. When two specialists are selected, a learned gate uses the initial physical history to combine their reconstructed fields.

\item \textbf{Multi-regime evaluation.}
We evaluate the framework on  incompressible-flow benchmarks spanning steady, Hopf-onset, and periodic regimes. The experiments assess autonomous specialist rollouts, cross-regime routing and fusion in regime overlaps, and the main local architectural choices, providing complementary evidence for the proposed local modeling and assembly strategy.

\end{itemize}




\par


REVIEWBLOCK00000


\par



\section{Related Work}
\paragraph{Projection, learned dynamics, and autonomous rollout.}
POD--Galerkin methods project the governing equations onto a reduced basis and retain an explicit reduced dynamical system \citep{sirovich1987,berkooz1993,benner2015}. Their long-horizon accuracy can deteriorate when truncated interactions are not adequately represented, motivating closure and stabilization strategies
\citep{amsallem2012}. Non-intrusive and data-driven ROMs instead infer reduced evolution laws directly from simulation data, ranging from projection-based operator inference to neural and latent-state dynamics
\citep{peherstorfer2016,hesthaven2018,champion2019,lee2020, pan2024domain,park2024}. Operator-inference methods provide a particularly direct connection between data-driven reduced dynamics and classical projection-based ROMs \citep{peherstorfer2016}. Neural-operator approaches learn mappings between function spaces rather than explicit low-dimensional state evolution\citep{lu2021,kovachki2023}. Hybrid ROMs retain projected or physics-based structure while learning missing or unresolved reduced dynamics \citep{duraisamy2019,fresca2022,dar2023}. RAL-MoE-ROM follows this hybrid principle, but places a structured sparse-MoE correction inside each regime-local ROM and trains the resulting specialist under multi-step  autonomous rollout.



\paragraph{Localized representations and local reduced dynamics.}
Localization has been explored through parameter-dependent reduced bases, localized operators and nonlinear manifolds \citep{amsallem2008,peherstorfer2014,kaiser2014,li2021cluster,diaz2024,bhat2025}. These approaches all reduce the burden on a single global representation, differing in whether localization is performed along parameter, state, or spatial coordinates. Recent work has further emphasized explicitly local reduced dynamics. More recently, ql-ROM partitions the solution manifold into clusters, constructs an intrusive ROM within each cluster, and connects the local models through assignment and change-of-basis operations
\citep{colanera2025}. Therefore, when complete local ROMs are constructed independently, their reduced coordinates and operators need not be directly compatible, so switching or combining them requires an explicit coordination mechanism. RAL-MoE-ROM takes a different approach: instead of coordinating local models in reduced-coordinate space, it lets local specialists evolve independently in their native coordinates and combines neighboring predictions only after reconstruction in the common physical space.





\paragraph{Mixture-of-experts and cross-chart assembly.}
Classical hierarchical mixtures learn conditional expert assignments \citep{jacobs1991,jordan1994}, while sparse MoE architectures activate only a subset of experts for each input \citep{shazeer2017,fedus2022}. MoE-based conditional specialization has also been introduced into scientific machine learning, including physics-informed constraints, 
neural-operator pretraining, and PDE forecasting
\citep{ICLR2024_9aeda582,NEURIPS2024_b83fae17,
NEURIPS2025_2d23a999,ICLR2026_754612bd}. These methods demonstrate the value of conditional specialization, but do
not address routing among complete autonomous ROMs. RAL-MoE-ROM therefore applies router  at two levels: sparse experts model unresolved dynamics within each local ROM, while an outer router acts on the complete regime-local specialists. The outer level selects a single specialist or an
applicable neighboring pair, with the latter combined through a history-conditioned gate after physical reconstruction.



\section{RAL-MoE-ROM}
\label{sec:method}

RAL-MoE-ROM separates reduced modeling into three complementary components: regime-local ROMs, autonomous local specialists, and cross-regime router and fusion. These components are organized by a two-level routing hierarchy. 

\subsection{Overlapping regime-local reduced charts}

Let $r\in\{S,H,P\}$ index steady, Hopf-onset, and periodic regimes, and let $\mathcal M_r$ denote the corresponding overlapping parameter regions. Each regime-local parameter region is equipped with its own reduced coordinate system. A physical velocity field is encoded and reconstructed as


REVIEWBLOCK00001


where $M_u$ is the finite-volume quadrature matrix. Gauge-corrected pressure is encoded analogously, yielding local coefficients
$\bm b^{(r)}$. 

A reduced chart comprises its mean fields, velocity and pressure bases, local coordinates, and normalization statistics. Each chart is associated with a regime-local parameter region and its own projected operators.
These quantities are constructed independently across regimes. Even when two charts have the same reduced dimension, their coordinates are not assumed to be componentwise aligned. Cross-regime assembly is consequently performed only after reconstruction in physical space.

\subsection{Regime-local hybrid specialists}

Within each chart, a projected Galerkin backbone describes the resolved dynamics, while a learned closure accounts for the unresolved reduced dynamics:


REVIEWBLOCK00002


The learned term $\bm\rho^u_{\theta_r}$ is the velocity closure associated with chart $r$. Its input $\bm\xi^{(r)}$ collects standardized local state variables, the projected right-hand side, the physical parameter, and a short history of preceding states.

Separate sparse MoE maps are used for the velocity correction $\bm\rho^u_{\theta_r}$ and the learned pressure correction $\bm\rho^p_{\theta_r}$:


REVIEWBLOCK00003


Routing is hierarchical within each specialist. An expert group $m^\star$ is selected by a chart-local router or predetermined for a single-group or fixed-group specialist. Within that group, the velocity and pressure channels $c\in\{u,p\}$ independently select their two highest-scoring routed experts, while a shared expert remains active. The nonnegative weights $\omega_e^c$ are normalized over the active experts. Each active expert combines a nonlinear feature-conditioned branch with explicit linear and low-rank quadratic responses of the local reduced state; the exact parameterization is given in Appendix~\ref{app:specialists}.

After advancing velocity over one reduced time step, pressure is reconstructed algebraically as


REVIEWBLOCK00004


Here $\mathcal Q_r^{PP}$ denotes the chart-local reduced Pressure--Poisson map, $\bm\rho^p_{\theta_r}$ provides the learned pressure correction, and $\bm\gamma_n^{(r)} =\operatorname{sigmoid}(\bm\eta^p_{\theta_r}(\bm\xi_n^{(r)}))$ modulates the Pressure--Poisson contribution componentwise; $\odot$ denotes componentwise multiplication. Thus velocity is advanced dynamically, whereas pressure is reconstructed algebraically rather than evolved as a separate recurrent state.


\subsection{Closed-loop rollout training}

Each specialist is initialized from an encoded physical history and is then advanced recursively using its own predicted states. To expose the learned correction to the states encountered during autonomous rollout, training unrolls each specialist for $K$ steps and minimizes the
normalized coefficient error


REVIEWBLOCK00005


Here $d_u^{(r)}$ and $d_p^{(r)}$ denote the chart-specific velocity and pressure coordinate dimensions, while $\bm\sigma_a^{(r)}$ and $\bm\sigma_b^{(r)}$ are scaling
vectors estimated from the corresponding regime-local training set; $\oslash$ denotes componentwise division. This coefficient-space objective is used for specialist training, whereas evaluation is performed on reconstructed physical fields. After initialization, no reference state from inside the prediction window is injected into the rollout. The chart-specific advancement rule is detailed in Appendix~\ref{app:specialists}, while rollout curricula and optimization settings are reported in Appendix~\ref{app:implementation}. Held-out trajectory splits are listed in Table~\ref{tab:chart_splits}.



\subsection{Cross-regime routing and physical-space T2-C assembly}

The outer E2 router receives only the standardized physical parameter $\widetilde\mu$ and returns a trajectory-level preference over the
regime-local specialists,
\[
\bm\pi(\mu)
=
\operatorname{softmax}\!\left(
\mathcal R_{E2}(\widetilde\mu)
\right).
\]
The probabilities are evaluated once and remain fixed throughout the rollout. The outer router does not use physical history. Chart-independent descriptors extracted from the initial history $\mathcal H_n$ enter only the T2-C gate described below.

A validation-supported applicability mask $m_r(\mathcal H_n,\mu,K)\in\{0,1\}$ marks whether specialist $r$ is applicable for the current parameter, initial history, and rollout horizon. The mask restricts the candidate set but does not modify the E2 probability
model itself. Under hard Top-1 routing, only the applicable specialist with the largest E2 probability is advanced.

When an adjacent pair is applicable ($S$--$H$ or $H$--$P$), both specialists are initialized from the same physical history, encoded in
their own coordinates, and advanced independently. Let $\bar\pi_r$ denote the E2 probability normalized over the selected pair. For reconstructed fields $\widehat{\bm q}^{(r)}$ and $\widehat{\bm q}^{(s)}$, T2-C predicts a history-conditioned logit correction to this pairwise preference and forms


REVIEWBLOCK00006


Here $\Delta_\psi=c_\psi([\widetilde\mu;\bm d_n])$ is a learned logit correction based on the standardized parameter and the chart-independent physical-history descriptor $\bm d_n$, and $\sigma$ denotes the sigmoid. The scalar $\alpha$ is computed once for the prediction window and is shared by velocity and pressure. The field vector $\bm q$ collects reconstructed velocity and gauge-corrected pressure on the common physical mesh.

The two local trajectories remain independent throughout the rollout: fusion is applied only to their reconstructed outputs and is never fed back into either specialist. If no specialist is applicable, the method abstains from producing a ROM prediction. Applicability is determined empirically from validation support and should not be interpreted as a formal stability guarantee.

\section{Experiments}
\label{sec:experiments}

We organize the experiments around three questions: whether the regime-local specialists support autonomous prediction, whether
physical-space fusion improves prediction near regime transitions, and how the main local architectural choices affect prediction.
The three incompressible-flow benchmarks play complementary roles in answering these questions. All are two-dimensional systems in which varying the Reynolds number traverses steady, Hopf-onset, and periodic regimes. The circular-cylinder benchmark provides the local architecture ablation, the centered-square benchmark evaluates cross-regime routing and fusion in the $S$--$H$ and $H$--$P$ overlaps, and the fluidic-pinball benchmark provides the autonomous comparison with a classical POD--Galerkin ROM. Table~\ref{tab:benchmark_summary} summarizes the benchmark scope and roles. All splits are made at the level of complete parameter trajectories rather than individual snapshots or rollout windows; chart-wise parameter ranges and split counts are reported in Appendix~\ref{app:implementation}.

We report relative physical-field errors $E_u$ and $E_p$ in percent; their weighted definitions and benchmark-specific aggregation rules are given in Appendix~\ref{app:implementation}. For routing and fusion comparisons, $E_{\mathrm{joint}}=E_u+E_p$. The sum of displayed component errors can differ slightly from the reported joint error because values are rounded independently. A rollout window is labeled ``Divergent'' if it violates the predefined rollout-divergence criterion, either through a non-finite state or a
threshold violation. N/E denotes a setting that is not evaluable under this criterion; table-specific notes identify the cause.

\par


REVIEWBLOCK00007


\par



\subsection{Do local specialists support accurate autonomous rollouts?}

We first examine whether the regime-local specialists maintain accurate predictions under autonomous rollout on the fluidic-pinball benchmark. Table~\ref{tab:pinball_autonomous} compares the learned specialists with a classical POD--Galerkin ROM over the Steady, Hopf, and Periodic regimes.

RAL-MoE-ROM achieves low physical-field errors across all four evaluation settings. In the Hopf and Periodic settings where POD--Galerkin errors are reported, the learned specialists consistently reduce both velocity and pressure
errors, with the largest differences appearing in pressure. For example, the Hopf joint error decreases from $48.94\%$ to $2.085\%$, while the standard Periodic error decreases from $43.20\%$ to $0.8415\%$. The longer-horizon Periodic-core evaluation shows the same trend, with $E_{\mathrm{joint}}$ reduced from $41.91\%$ to $1.672\%$. Together, these results show that the learned specialists maintain low physical-field error over the evaluated autonomous horizons.

\par


REVIEWBLOCK00008


\par

A complementary POD--DeepONet experiment further separates one-step prediction from autonomous rollout.
The model attains $E_u=1.5921\%$ and $E_p=23.3447\%$ on 4,228 held-out one-step pairs, but only one of 28 held-out trajectories completes the prescribed autonomous rollout evaluation. Because its temporal cadence differs from that of the Periodic specialist, we use this result only to contrast one-step prediction with autonomous rollout, rather than as a matched long-horizon accuracy comparison (Appendix~\ref{app:pinball}).





\subsection{Does physical-space fusion improve prediction in regime overlaps?}
To evaluate the second routing level while holding the local specialists fixed, we compare the centered-square $S$--$H$ and $H$--$P$ overlaps at $K=24$ (Table~\ref{tab:routing_fusion}). We compare hard E2 Top-1 routing with two simple fusion rules, the Equal blend and the E2-probability blend, together with the learned T2-C gate; a diagnostic convex oracle provides an a posteriori reference. To test whether T2-C benefits specifically from physical-history information, we additionally include an architecture-matched T2-C $\mu$-only ablation. The same gate architecture and optimization protocol are retained, but the standardized history-descriptor channels are set to zero. This preserves the gate parameterization and trainable parameter count while removing history information.


The $\mu$-only ablation reveals that parameter conditioning alone has overlap-dependent value. For $S$--$H$, its joint error is essentially unchanged relative to the E2-probability blend ($2.1846\%$ versus $2.1847\%$), whereas for $H$--$P$ it improves the same baseline from $1.6618\%$ to $1.1504\%$. Adding the physical-history descriptor improves both overlaps further, reducing $E_{\mathrm{joint}}$ to $1.1958\%$ and $0.8465\%$, respectively. Relative to the a posteriori best fixed-specialist reference, full T2-C reduces joint error by $30.9\%$ on $S$--$H$ and $29.9\%$ on $H$--$P$. Relative to the deployable E2 Top-1 router, the corresponding reductions are $30.9\%$ and $54.2\%$. The validation-selected T2-C errors are also close to the diagnostic convex oracle in both overlaps. Representative reconstructed fields and velocity-error maps are shown in Figure~\ref{fig:square_overlap}.

Across three paired gate-training seeds, full T2-C outperforms the architecture-matched $\mu$-only gate in every run for both overlaps.
Across the three paired runs, the mean reduction relative to $\mu$-only is $45.3\%$ for $S$--$H$ and $26.2\%$ for $H$--$P$. Full T2-C yields
$E_{\mathrm{joint}}=1.1957\pm0.0001\%$ on $S$--$H$ and $0.8500\pm0.0035\%$ on $H$--$P$, where $\pm$ denotes the sample standard
deviation across the three gate-training seeds. Only gate training is repeated across seeds; the specialists, E2 router, candidate rollouts, data splits, normalization, and optimization protocol are fixed. These statistics therefore characterize gate-training variability rather than full-pipeline uncertainty. Component-wise errors and seed-level results are reported in Appendix~\ref{app:square}, with the corresponding fluidic-pinball overlap study in Appendix~\ref{app:pinball}.


\par


REVIEWBLOCK00009


\par

\par


REVIEWBLOCK00010


\par



\subsection{How do local architectural choices affect prediction?}

The circular-cylinder benchmark is used to examine three local architectural choices: structured expert parameterization, the projected ROM backbone, and regime localization (Table~\ref{tab:circular_full_reorg}). Vanilla-FNN-MoE retains the regime-local setting but replaces the structured expert map with conventional feed-forward experts. DataOnly-MoE removes the projected Galerkin and Pressure--Poisson backbone, leaving the learned local dynamics to account for the evolution from data. Global MoE replaces the regime-local charts and specialists with a single global reduced representation and dynamical model. All comparisons are reported in physical-field error so that models with different reduced coordinates are assessed in a common physical space. These comparisons characterize architectural effects under the available fixed models rather than parameter-matched capacity controls. Total and active trainable parameter counts for the analyzed Hopf and Periodic specialists are reported in Appendix~\ref{app:parameter_counts}.



REVIEWBLOCK00011


















Replacing the structured expert map with conventional feed-forward experts increases both velocity and pressure errors in the Steady and Periodic regimes; the Vanilla-FNN-MoE Hopf model is not evaluable under the prescribed rollout criterion. Removing the projected ROM backbone produces the largest degradation, with particularly large pressure errors for DataOnly-MoE in the Steady and Hopf regimes. Among the ablations, Global MoE provides the closest overall comparison and is particularly close to the full model in the Hopf regime. Nevertheless, the regime-local specialists retain lower reported physical-field errors in all three regimes. Together, these comparisons support the roles of structured expert corrections, the projected ROM backbone, and regime-local representations under the tested architectures. They are not parameter-matched capacity controls. Cross-regime assembly is evaluated separately in Table~\ref{tab:routing_fusion}. Additional post-hoc internal-routing diagnostics are reported in
Appendix~\ref{app:internal_routing}.




\FloatBarrier
\section{Conclusion}

We introduced RAL-MoE-ROM, a hierarchical reduced-order modeling framework that separates within-regime dynamical specialization from cross-regime assembly. Regime-local ROMs retain independent reduced coordinates and augment projected dynamics with structured sparse corrections, while
cross-regime coupling is performed only after reconstruction in physical space. Across three parameterized incompressible-flow benchmarks, the local specialists support accurate autonomous rollouts, and history-conditioned T2-C improves prediction in the evaluated regime overlaps. The present study assumes validation-supported regime coverage and evaluates pairwise fusion only for adjacent specialists on two-dimensional incompressible flows. Future work will investigate learned regime discovery, constraint-preserving physical-space assembly, and higher-dimensional dynamical systems.


\subsubsection*{Reproducibility statement}
The main text reports the benchmark roles, evaluation horizons, baselines, and metrics. Appendices~A--E document the local-coordinate construction, projected operators, specialist architecture, routing and fusion procedure, and implementation details. Appendices~F--H provide benchmark-specific results.
Appendix~\ref{app:internal_routing} reports post-hoc internal-routing statistics
from fixed validation-selected Hopf and Periodic specialists on held-out
rollout windows; these diagnostics are not used for model selection.
Appendix~\ref{app:square} additionally reports an architecture-matched descriptor
ablation and three-seed variability of the centered-square T2-C gate, with the
underlying specialists and E2 router held fixed.

\subsection*{AI use statement}
% Authors must confirm this disclosure against the final submission workflow.
Generative AI tools, including ChatGPT and Codex, were used solely to assist with language polishing and LaTeX formatting. The authors take full responsibility for reviewing and verifying all scientific
content, mathematical formulations, experimental configurations, numerical
results, and conclusions, and for the final manuscript and all reported results.

\bibliography{references,legacy_references}
\bibliographystyle{iclr2027_conference}

\appendix
\section{Full problem setup and weighted POD construction}
\label{app:problem_pod}

\subsection{Parameterized full-order model, gauge, and regime coverage}

After spatial discretization, the velocity
$\bm u(t;\mu)\in\mathbb R^{N_u}$ and pressure
$p(t;\mu)\in\mathbb R^{N_p}$ satisfy


REVIEWBLOCK00012


Pressure is defined only up to an additive constant.
With positive finite-volume weights $\bm w$, each pressure snapshot is mapped
to the weighted zero-mean gauge


REVIEWBLOCK00013



The parameter domain is covered by overlapping regime-local parameter regions,


REVIEWBLOCK00014


where $S$, $H$, and $P$ denote the steady, Hopf-onset, and periodic
regimes, respectively.
Only adjacent overlaps are used; no direct $S$--$P$ overlap is considered.


\subsection{Area-weighted local POD}

The finite-volume quadrature induces the weighted inner products


REVIEWBLOCK00015



For chart $r$, let $d_u^{(r)}$ and $d_p^{(r)}$ denote the retained velocity
and pressure dimensions.
Let $X_u^{(r)}$ and $X_p^{(r)}$ collect the centered velocity and
gauge-corrected pressure snapshots from the training trajectories assigned
to that chart.
The corresponding chart-specific means are denoted by
$\bar{\bm u}^{(r)}$ and $\bar p^{(r),\circ}$.
The local POD bases solve


REVIEWBLOCK00016


where


REVIEWBLOCK00017



Equivalently, consider the weighted singular value decompositions


REVIEWBLOCK00018


The retained bases are


REVIEWBLOCK00019


and satisfy


REVIEWBLOCK00020


Because the pressure snapshots and their chart mean satisfy the same weighted
gauge condition, the pressure basis also satisfies


REVIEWBLOCK00021



Encoding and decoding are performed independently within each chart:


REVIEWBLOCK00022



The reduced coefficients are standardized using chart-local statistics
estimated from the corresponding training trajectories:


REVIEWBLOCK00023


Additional specialist features and their normalization are specified in
Appendix~\ref{app:specialists}.

All chart-local means, POD bases, and scaling statistics are constructed
independently.
Consequently, reduced coefficient vectors from different charts are not
assumed to be componentwise comparable; only their decoded physical fields lie in a common physical comparison space.






\section{Projected Galerkin and pressure operators}
\label{app:rom}

\subsection{Chart-local Galerkin momentum dynamics}

For each chart $r\in\{S,H,P\}$, substituting the chart-local affine reconstructions
from Appendix~\ref{app:problem_pod} into the semi-discrete
momentum equation and testing against the local velocity basis yields


REVIEWBLOCK00024


Here ``$:$'' denotes the double contraction.
For velocity modes
$\bm\phi_i^{u,(r)}$,
$i=1,\ldots,d_u^{(r)}$,
and pressure modes
$\phi_\ell^{p,(r)}$,
$\ell=1,\ldots,d_p^{(r)}$,
the projected operators are


REVIEWBLOCK00025



We denote the resulting chart-local projected momentum vector field by


REVIEWBLOCK00026


This vector field provides the explicit projected Galerkin backbone of the
regime-local specialist introduced in
Eq.~(\plaineqref{eq:main_specialist}).


\subsection{Gauge-consistent pressure--Poisson map}

The pressure component is represented through a chart-local
Pressure--Poisson relation.
At the continuous level, the gauge-corrected pressure satisfies


REVIEWBLOCK00027


with boundary conditions inherited from the momentum equation.
After discretization, substitution of the chart-local affine velocity
representation, and projection onto the local pressure basis, the reduced
pressure coefficients satisfy


REVIEWBLOCK00028


The corresponding chart-local Pressure--Poisson map is therefore


REVIEWBLOCK00029



Because the pressure basis lies in the weighted zero-mean subspace defined in
Appendix~\ref{app:problem_pod}, the constant-pressure null direction is
removed before projection.
Thus $\mathcal Q_r^{PP}$ produces pressure coefficients in the native
coordinates of chart $r$ without introducing an additional cross-chart
pressure gauge.

The pair


REVIEWBLOCK00030


is constructed independently for each regime-local chart.
Accordingly, each chart carries its own parameter region, mean fields,
POD bases, coefficient scalings, projected momentum operators, and
Pressure--Poisson map.
These chart-local operators are used only within their native reduced
coordinates; no componentwise interpolation of reduced states or projected
operators is assumed across charts.


\section{Regime-local specialist architecture and closed-loop training}
\label{app:specialists}

For each regime-local chart $r$, we define a self-contained 
specialist in its native reduced coordinates: velocity is advanced
dynamically, whereas pressure is reconstructed algebraically.
The projected operators from Appendix~\ref{app:rom} provide the underlying physics, while learned terms model unresolved reduced dynamics. The velocity coefficients satisfy


REVIEWBLOCK00031


After each reduced velocity time step, pressure is reconstructed through


REVIEWBLOCK00032


The gate modulates the projected Pressure--Poisson contribution, while
$\bm\rho^p_{\theta_r}$ provides the learned pressure correction.


\subsection{Features and structured sparse experts}

Let


REVIEWBLOCK00033


denote the projected velocity right-hand side defined in Eq.~(\plaineqref{eq:app_galerkin_backbone}). The specialist uses the ordered three-state history


REVIEWBLOCK00034


and constructs


REVIEWBLOCK00035


The feature map contains standardized local states, projected
right-hand sides and parameter features. All normalization statistics are estimated from the corresponding chart-specific training set. Internal routing first selects an expert group. When group selection is learned,


REVIEWBLOCK00036


whereas $m^\star$ is predetermined for single-group or fixed-group
specialists. Within the selected group, velocity and pressure use separate channel routers. For $c\in\{u,p\}$, let $\mathcal T_2^c$ denote the two routed experts with the largest channel-specific weights. Together with the shared expert $e=0$,


REVIEWBLOCK00037


The two channel mixtures define the velocity and pressure corrections
$\bm\rho_{\theta_r}^u$ and $\bm\rho_{\theta_r}^p$, respectively. The explicit linear and quadratic branches operate on the channel-specific reduced-state inputs


REVIEWBLOCK00038


For output component $j$, each expert is


REVIEWBLOCK00039


Thus each expert combines a nonlinear feature-conditioned response with
explicit linear and low-rank quadratic responses of the local reduced state.


\subsection{Frozen-context advancement and closed-loop training}

During one reduced velocity time step, the current pressure, recent history, and
physical parameter define a context $\mathcal C_n^{(r)}$ that is held fixed
while the velocity state is advanced. We write


REVIEWBLOCK00040


The chart-specific numerical update is denoted by


REVIEWBLOCK00041


Continuous-time specialist instances use Runge--Kutta integration with the
same frozen context across stage evaluations; benchmark-specific stepping
details are reported in Appendix~\ref{app:implementation}.
The updated velocity is then supplied to
Eq.~(\plaineqref{eq:app_pressure_closure}), after which the history is
shifted and autonomous rollout continues.

For a closed-loop rollout of length $K$, the specialist recursively
consumes its own predicted states and minimizes


REVIEWBLOCK00042


After initialization, no reference state inside the prediction window is
injected into the autonomous rollout. Training uses progressively longer
chart-specific rollout horizons; the corresponding curricula and
optimization settings are listed in Appendix~\ref{app:implementation}.

\section{Outer routing and T2-C details}
\label{app:routing}

\paragraph{Parameter-only regime routing.}
The E2 router receives only the standardized scalar parameter,


REVIEWBLOCK00043


The trajectory-level preference is evaluated once and remains fixed
throughout the subsequent rollout. The regime graph contains only adjacent transitions,


REVIEWBLOCK00044


so direct $S$--$P$ fusion is excluded.
Let
$m_j(\mathcal H_n,\mu,K)\in\{0,1\}$
denote the validation-supported applicability of specialist $j$ for the
given initialization and rollout horizon.
The E2 preference restricted to applicable specialists is


REVIEWBLOCK00045


A single applicable specialist yields the Top-1 prediction.
Top-2 fusion is considered only for an applicable pair belonging to
$\mathcal E_{\rm adj}$.
If no specialist is applicable, the hierarchy abstains from prediction.

For an ordered selected pair $(r_1,r_2)$, the pair-normalized E2
preference is


REVIEWBLOCK00046




\paragraph{Chart-independent physical-history descriptor.}
E2 is intentionally parameter only.
T2-C augments this trajectory-level preference with a chart-independent
descriptor computed from the current physical state and the two preceding
states.
The descriptor summarizes field magnitude, first-order temporal variation,
velocity growth, and second-order velocity variation.

Let
$|\Omega|=\bm 1^\top\bm w$,
$\Delta t_n=t_n-t_{n-1}$,
and let $U_{\rm ref}$, $P_{\rm ref}$, and $T_{\rm ref}$ denote fixed
reference scales.
For gauge-corrected pressure $p_n^\circ$, define the normalized field
magnitudes


REVIEWBLOCK00047


and the corresponding first-difference rates


REVIEWBLOCK00048


A velocity-growth indicator is


REVIEWBLOCK00049


To capture changes in the local temporal trend, define


REVIEWBLOCK00050



The resulting descriptor is


REVIEWBLOCK00051


Because $\bm d_n$ is computed entirely from reconstructed physical fields,
it is independent of either specialist's native POD coordinates.


\paragraph{Physical-space T2-C fusion.}
The logit-correction network receives
$[\widetilde\mu;\bm d_n]$ and uses two hidden layers of width 64 with SiLU
activations followed by a scalar linear output.
The final layer is initialized at zero, so the initial correction vanishes
and the fusion weight reduces to the E2 pair preference.
Specifically,


REVIEWBLOCK00052



The two selected specialists encode the same physical history in their
respective local charts and then evolve independently.
The scalar $\alpha$ is fixed throughout the prediction window and is shared
by velocity and pressure.
After reconstruction on the common physical mesh, T2-C forms


REVIEWBLOCK00053


where the pressure field is gauge corrected.
The fused field is an output only and is never projected back into either
specialist.

The correction network is trained directly on the reconstructed physical
fields. For rollout windows $w$ and prediction steps $k$,


REVIEWBLOCK00054


Optimization and checkpoint-selection details are reported in
Appendix~\ref{app:implementation}.

\section{Reproducibility and implementation details}
\label{app:implementation}

\paragraph{Physical-field error metrics.}
We evaluate reconstructed velocity and gauge-corrected pressure using the relative weighted field errors


REVIEWBLOCK00055


Here $\|\bm v\|_M=(\bm v^\top M\bm v)^{1/2}$ uses the finite-volume
quadrature introduced in Appendix~\ref{app:problem_pod}; the velocity norm includes both components, and $p^\circ$ denotes pressure in the common zero-mean gauge. Unless stated otherwise, benchmark-specific temporal and trajectory aggregation is reported with the corresponding experiment. Whenever a joint score is reported,
$E_{\mathrm{joint}}=E_u+E_p$. Independently rounded component values may therefore differ slightly from the displayed joint score.


\paragraph{Trajectory-level partitioning.}
Table~\ref{tab:chart_splits} summarizes the resulting chart-wise parameter coverage, reduced dimensions, trajectory splits, and evaluation horizons. All splits are defined at the level of complete parameter trajectories rather than individual snapshots or rollout windows. Local bases, normalization statistics, and model parameters are constructed from training trajectories only; validation trajectories are used for model and checkpoint selection, and test trajectories are reserved for the reported evaluations. 

\par


REVIEWBLOCK00056


\par


\paragraph{Centered-square T2-C gate optimization.}
Both full T2-C and the architecture-matched $\mu$-only gate use AdamW
(learning rate $3\times10^{-4}$, weight decay $10^{-4}$), 8,000 steps,
and gradient-norm clipping at 1. Minibatches sample training windows with
replacement. Validation is evaluated at step 1 and every 100 steps;
checkpoint selection minimizes the largest Reynolds-number-wise full-window mean joint
error, breaking ties by the pooled mean, without test-based selection.
The paired-seed evaluation is reported in Appendix~\ref{app:square}.



\paragraph{Autonomous rollout protocol.}
All reported specialist evaluations are autonomous after initialization from the prescribed physical history: no reference state from inside the prediction window is injected during rollout.
Each specialist is advanced using its chart-specific time-advancement rule described in Appendix~\ref{app:specialists}. For continuous-time instances, Runge--Kutta stage evaluations share the same frozen pressure/history context within one reduced velocity time step. After each reduced time step, the pressure reconstruction and local history are updated before the next autonomous prediction.


\paragraph{Benchmark discretization and sampling.}
The centered-square dataset contains 120 complete trajectories over
$Re\in[50,150]$ on a 9,400-cell finite-volume mesh. The CFD time step is $\delta t_{\rm CFD}=0.02$, and snapshots are stored every 200 CFD steps, giving $\Delta t_{\rm snap}=4$. The fluidic-pinball dataset contains 131 complete trajectories over $Re\in[10,60]$ on a 70,194-cell mesh, with $\delta t_{\rm CFD}=0.005$ and a snapshot interval of $0.25$ physical time units. The circular-cylinder dataset contains 100 Reynolds-parameterized trajectories spanning
$Re\in[20,200]$ on a common finite-volume mesh represented by 97,368 spatial points. Trajectories contain different numbers of snapshots, yielding 14,167 snapshots in total. The unsteady simulations employ adaptive CFD time stepping subject to $Co_{\max}=0.8$ and $\delta t_{\max}=0.2$, while the snapshot interval is adjusted according to the estimated vortex-shedding period, with eight snapshots retained per shedding cycle. The chart-wise parameter ranges, reduced dimensions, splits, and evaluation horizons are summarized in Table~\ref{tab:chart_splits}.


\paragraph{Specialist optimization.}
Table~\ref{tab:specialist_training} summarizes the reported training configurations, which use AdamW with chart-specific learning rates, weight decay, and rollout curricula.
Additional batching, scheduler, and precision settings are provided in the supplementary implementation.

\par


REVIEWBLOCK00057


\par

\section{Circular-cylinder results}
\label{app:circular}

Table~\ref{tab:circular_full} reports the min--max physical-field error across held-out parameter trajectories under autonomous rollout. In the Hopf regime, the Vanilla-FNN-MoE model does not satisfy the predefined validation-stage rollout criterion; no held-out result is therefore reported.

\par


REVIEWBLOCK00058


\par

In the $S$--$H$ overlap, E2 Top-1 yields $E_u=0.0357$--$0.0410\%$ and
$E_p=1.6746$--$1.9028\%$ across the held-out trajectories. T2-C lowers these ranges to $0.0148$--$0.0355\%$ and $0.7387$--$1.7212\%$, respectively. The corresponding mean weight assigned to the Steady specialist varies from $0.0861$ to $0.9987$, indicating that the learned assembly is not equivalent to a fixed blend across the overlap. For reference, the Global MoE gives $E_u=0.0851$--$0.0953\%$ and $E_p=1.9517$--$2.3106\%$ under the same overlap evaluation.

\FloatBarrier
\section{Centered-square results}
\label{app:square}


This appendix provides the complete centered-square overlap results
underlying the main-text routing and fusion comparison. Table~\ref{tab:square_fusion_full} reports the velocity, pressure,
and joint errors for the individual specialists and the fusion baselines in the $S$--$H$ and $H$--$P$ overlap regions. In the main text, ``Best fixed'' denotes the fixed specialist with the
lowest aggregate joint error over the corresponding overlap, whereas the convex oracle is a window-wise a posteriori reference.



\paragraph{Overlap evaluation and descriptor ablation.}
For both overlaps, $E_u$, $E_p$, and $E_{\mathrm{joint}}=E_u+E_p$ are reported at the terminal rollout step ($k=24$) and averaged over 16 held-out windows.  This terminal metric differs from the full-window objective used by the diagnostic convex oracle. To assess the role of history information in the second routing level, we also include T2-C $\mu$-only, which uses the same correction-gate architecture and training protocol as T2-C but removes the standardized history descriptor by setting its channels to zero. Table~\ref{tab:square_fusion_full} reports the validation-selected comparison, and the corresponding three-seed gate variability is summarized separately below.

\par


REVIEWBLOCK00059


\par

\par


REVIEWBLOCK00060


\par

The regime-local specialists also exhibit low aggregate physical-field errors under their respective held-out evaluations: $(E_u,E_p)=(1.5311,1.4583)\%$ in the Steady regime, $(0.4387,0.4844)\%$ in the Hopf regime, and $(0.9395,1.0921)\%$ in the Periodic regime. Within the overlap evaluations, however, the more accurate individual specialist differs by interface: the Hopf specialist is the better single-model reference in $S$--$H$, whereas the Periodic specialist is more accurate in $H$--$P$. This overlap dependence motivates an adaptive assembly rule rather than a fixed specialist choice across neighboring regimes.





\paragraph{Gate-training seed variability.}
To assess the sensitivity of the second routing level, we retrain only the
T2-C gate for seeds 42001, 42002, and 42003, keeping the specialists, E2
router, candidate rollouts, data splits, and normalization fixed. Full T2-C
and the architecture-matched $\mu$-only variant share the optimization and
selection protocol in Appendix~\ref{app:implementation}.
Table~\ref{tab:t2c_seed_variability} reports the mean and sample standard
deviation ($N=3$, denominator $N-1$), and Table~\ref{tab:t2c_seedwise}
lists the seed-wise joint errors. These statistics describe gate-training
variability, not across-$Re$ variation, a confidence interval, or
full-pipeline uncertainty.

\par


REVIEWBLOCK00061


\par

\par


REVIEWBLOCK00062


\par

Full T2-C improves upon T2-C $\mu$-only for all three paired seeds in both overlaps. The corresponding reductions in mean joint error are $45.28\%$ in $S$--$H$ and $26.16\%$ in $H$--$P$. In the $S$--$H$ overlap, the $\mu$-only gate remains close to the E2-probability blend, whereas in $H$--$P$ it provides a moderate correction.
In both cases, incorporating the history descriptor yields a further and consistent improvement. These results support the use of history-conditioned assembly in the present centered-square setting.


Figures~\ref{fig:app_square_sh_full} and \ref{fig:app_square_hp_full} provide complete views of the velocity, gauge-corrected pressure, and pointwise error. In the S--H overlap, the Steady and Hopf predictions differ mainly in the near-wake recirculation region and in the downstream pressure-recovery pattern. The T2-C reconstruction remains visually close to the more accurate Hopf prediction while reducing localized discrepancies in both velocity and pressure, consistent with the error reductions reported in Tables~\ref{tab:square_fusion_full} and \ref{tab:square_parameter}.

In the H--P overlap, the discrepancy between the Hopf and Periodic specialists is more pronounced around the obstacle corners, the separated shear layers, and the developing wake. The T2-C reconstruction reduces these localized errors while preserving the dominant large-scale wake structure, consistent with the substantially lower joint velocity--pressure errors reported for this overlap.

\par


REVIEWBLOCK00063


\par

\par


REVIEWBLOCK00064


\par

\FloatBarrier
\section{Fluidic-pinball results}
\label{app:pinball}

The POD--DeepONet baseline provides a complementary comparison between
one-step prediction and autonomous rollout. It attains $E_u=1.5921\%$ and $E_p=23.3447\%$ on 4,228 held-out input--output pairs, but only 1 of 28 held-out trajectories completes the prescribed autonomous rollout evaluation. This difference illustrates that one-step accuracy does not by itself imply reliable autonomous prediction. Because the POD--DeepONet evaluation uses a different temporal cadence from
the Periodic specialist, it is used only for this comparison of one-step prediction and autonomous rollout and not as a cadence-matched long-horizon accuracy baseline.

\paragraph{Periodic-core evaluation protocol.}
The $K=32$ Periodic result in Table~\ref{tab:pinball_autonomous} corresponds to the standard periodic evaluation. We additionally report a longer $K=56$ evaluation restricted to the Periodic-core final-test trajectories. The full-order snapshots are stored every $0.25$ physical time units; subsampling every fourth snapshot gives an effective interval $\Delta t=1$, matching the Periodic specialist's native history cadence. The $K=56$ setting therefore provides a longer autonomous evaluation within the same periodic regime without changing the specialist representation or
sampling convention.

\par


REVIEWBLOCK00065


\par

T2-C improves upon Candidate~1 in all four overlap evaluations.
The relative reductions in joint error are approximately $1.9\%$ and
$2.3\%$ for $S$--$H$ at $K=24$ and $K=56$, and $4.5\%$ and $4.9\%$
for $H$--$P$, respectively. Although these gains are more modest than those observed for the centered-square benchmark, they persist at both rollout horizons. T2-C also remains close to the corresponding window-wise convex oracle, indicating that the learned fusion captures most of the available benefit from convex physical-space assembly in these overlap evaluations.

\FloatBarrier
\section{Additional internal-routing diagnostics}
\label{app:internal_routing}

\paragraph{Evaluation protocol and routing statistics.}
We analyze the frozen circular-cylinder Hopf and Periodic specialists on their held-out autonomous rollout windows.
For the Periodic specialist, we use 13 fixed windows at each of four held-out Reynolds numbers, giving 52 windows of length $K=48$. For the Hopf specialist, we use five fixed windows at each of three held-out Reynolds numbers, giving 15 windows of length $K=56$. One routing record is collected per reduced time step at the first integration stage, resulting in 2496 Periodic and 840 Hopf routing decisions per channel. Thus the reported statistics characterize reduced time step routing rather than all intermediate right-hand-side evaluations. Overlapping windows are not interpreted as independent trajectories, and these diagnostics are not used for model selection.

For $N$ routing decisions, let $I_{n,e}\in\{0,1\}$ indicate whether routed expert $e$ is selected at decision $n$, with the group identity included in the expert label. We define the normalized activation frequency


REVIEWBLOCK00066


where the factor $2$ accounts for the Top-2 routed experts selected at each decision, so that $\sum_e f_e=1$. The corresponding mean routing mass is


REVIEWBLOCK00067


where $\omega_{n,e}$ denotes the final routed mixing coefficient assigned to expert $e$. Experts that are never selected are retained with $f_e=m_e=0$.

We summarize the dispersion of expert selection by the normalized utilization entropy


REVIEWBLOCK00068


where $N_E$ is the number of routed experts available to the corresponding channel. Because $H$ is computed from activation frequency rather than routing mass, it characterizes diversity of expert selection rather than equality of the assigned weights. Distinct Top-2 pairs retain group identity but are treated as unordered.

For the two specialists considered here, the architecture assigns a fixed coefficient $4/7$ to the shared expert and a total coefficient $3/7$ to the routed Top-2 component. Consequently,


REVIEWBLOCK00069


The routing mass therefore measures mixture allocation; it does not quantify the norm of an expert output or, by itself, establish a physical interpretation.
Labels such as \texttt{g0:e0} refer only to routed experts and should not be confused with the shared expert $e=0$ in
Eq.~(\plaineqref{eq:app_channel_sparse_moe}).

\par


REVIEWBLOCK00070


\par

\paragraph{Periodic routing profile.}
The Periodic specialist uses a broad range of its available routed experts. Across the held-out rollouts, the velocity channel selects all 18 routed experts and realizes 33 distinct group--pair combinations, while the pressure channel selects 14 of 18 experts through 16 combinations (Table~\ref{tab:internal_routing_summary}). The most frequent pair accounts for only $16.27\%$ of velocity decisions and $31.25\%$ of pressure decisions, consistent with the higher utilization entropy of the velocity channel. The selected groups also vary across the four held-out Reynolds numbers.
Figure~\ref{fig:periodic_internal_routing} summarizes the corresponding group frequencies and routed-expert masses.



\paragraph{Hopf routing profile.}
The Hopf specialist exhibits a substantially more concentrated routing pattern. Because it contains a single expert group, only channel-level routing is active. Across all evaluated windows, the velocity channel selects the pair
$(e_1,e_3)$ and the pressure channel selects $(e_2,e_4)$ at every recorded reduced time step. A separate stage-wise check confirms that these same pairs are retained across all four Runge--Kutta stages.

Although the selected support is fixed, the routed weights are strongly nonuniform. Velocity expert $e_3$ receives approximately $99.93\%$ of the total routed mass, while pressure expert $e_2$ receives $84.16\%$. The secondary pressure expert $e_4$ nevertheless receives a Reynolds-number-dependent weight, with its mean mass increasing across the three held-out parameter values.

The two specialists therefore exhibit different internal utilization patterns. The Periodic specialist distributes routing decisions across many group--pair combinations, whereas the Hopf specialist retains fixed sparse support and adapts primarily through the mixture weights. The single Hopf group is part of the prescribed architecture and should not be interpreted as collapse of a learned multi-group router.

\paragraph{Periodic rollout-position diagnostic.}
Figure~\ref{fig:periodic_routing_step} shows the mean routed-expert mass as a function of normalized rollout position, using the same 52 held-out Periodic windows. The horizontal coordinate reflects prediction position within the
autonomous window rather than physical oscillation phase. This diagnostic complements the aggregate statistics above by showing how expert allocation varies over the autonomous rollout.


\par


REVIEWBLOCK00071


\par

\par


REVIEWBLOCK00072


\par

\clearpage
\section{Trainable and active parameter counts}
\label{app:parameter_counts}

\paragraph{Counting convention.}
Total parameters count unique trainable parameters in each analyzed
circular-cylinder specialist. Active parameters count unique trainable
parameters used by one batch-size-one prediction-path routing decision,
not their union across a batch, trajectory, or all RK4 stages. These include
the encoder and refinement blocks, executed routers, the selected group's
shared experts, channel-specific Top-2 velocity and pressure experts, and
the learned pressure correction. Fixed POD bases, projected operators,
normalization statistics, and non-trainable buffers are excluded.

\par


REVIEWBLOCK00073


\par

The Hopf and Periodic specialists activate 44.53\% and 17.09\% of their
trainable parameters per prediction decision, respectively. These counts
describe conditional parameter activation, not measured FLOPs, runtime,
or speedup. Complete parameter-matched counts are unavailable for the
Vanilla-FNN-MoE, DataOnly-MoE, and Global MoE ablations, so the local
architecture comparison is not interpreted as a capacity-matched study.


\end{document}
